\documentclass[10pt]{article}
\usepackage{icmlko}

\icmlmeta{IP 파운데이션 월드모델}{252}{천장에 닿아 있다}{축 열일곱으로 남은 여지는 $+$0.013 이고, 아홉 중 다섯에서 풀링이 이미 도메인 전용 모형을 이긴다}

\begin{document}

\twocolumn[
\icmlheader

\begin{abstract}
노트 240--251은 열두 편 내리 \emph{측정}에 관한 것이었고 판은 0.4569에서
안 움직였다. 방향을 바꿔 물었다 --- \textbf{지금 축으로 도달 가능한 천장은
얼마인가.} 유보에 직접 적합하면 판이 0.5465 로 보이지만 그것은 거짓이다.
\textbf{라벨을 섞어 같은 적합을 하면 0.2080 이 나온다} --- 과적합 바닥이
그만큼이다(아이돌은 섞은 라벨로도 0.68, 팝업 0.55). 정직한 눈금은 유보
\emph{안}에서의 5겹 CV 이고, 그것이 \textbf{0.4697} 이다. \textbf{남은 여지는
$+$0.0128 뿐이다.} 그리고 도메인별로 보면 더 세다 --- \textbf{아홉 중 다섯에서
여지가 음수다.} 게임 $-$0.039 · 펀딩 $-$0.046 · 세계애니 $-$0.033 · 모바일
$-$0.026 · 애니 $-$0.014. \textbf{풀링 모형이 그 도메인의 유보 라벨을 직접
보고 배운 전용 모형보다 낫다.} 풀링은 타협이 아니라 이득이다. 여지가 진짜로
있는 곳은 \textbf{도서 하나}다($+$0.177, 몫 6.2\%). 왜인지는 라벨을 안 보고도
보인다 --- \textbf{도서는 학습 80건에 유보 163건으로, 학습이 유보보다 적은
유일한 도메인}이고(다른 곳은 1.4$\sim$8.8배) 학습이 1999--2025에 흩어져 있고
축 넷이 덮음 0\%다. 즉 도서의 여지는 모형이 아니라 \textbf{수집}의 문제다.
비($=$학습/유보)로 여지를 예측해 봤는데 spearman $-$0.52(n$=$9)로 법칙이라
부를 수 없다 --- 도서가 뚜렷한 이상치일 뿐이다. \textbf{프로그램 수준의 결론}:
지금 축은 거의 소진됐다. 다음 이득은 더 나은 모형이 아니라 \textbf{새 정보}에서
온다.
\end{abstract}
]

\section{열두 편 뒤에 방향을 바꿨다}

노트 240부터 251까지 열두 편이 전부 \emph{어떻게 재는가}였다 --- 겹말 ·
안쪽 눈금 · 상쇄 · 교차항. 배운 것은 많지만 판은 노트 242의 0.4569에서
한 번도 안 움직였다. 그래서 물음을 바꿨다: \textbf{움직일 여지가 남아
있기는 한가.}

값싼 눈금이 있다. 유보에 \emph{직접} 적합해 보면, 지금 축으로 그 표본을
설명할 수 있는 한계가 나온다. 순위에는 안 쓴다 --- 노트 222의 도메인
천장표처럼 어디에 여지가 있는지 보는 자다.

\section{첫 수는 거짓이었다}

\begin{center}\small
\begin{tabular}{lrr}
\hline
 & 판 & \\
\hline
지금(풀링, 배포) & 0.4569 & \\
유보에 직접 적합(선형) & 0.5465 & $+$0.0895 \\
유보에 직접 적합(나무) & 0.9\ldots & \\
\hline
\end{tabular}
\end{center}

나무가 0.9를 넘는 것이 단서였다. 34개 열로 25$\sim$700행을 맞히니 외우는
것이다. 그래서 \textbf{라벨을 섞어 같은 적합}을 했다 --- 신호가 0 인데도
나오는 값이 과적합 바닥이다.

\begin{center}\small
\begin{tabular}{lrrr}
\hline
도메인 & $n$ & 속천장 & \textbf{섞음} \\
\hline
아이돌 & 25 & 0.797 & \textbf{0.682} \\
팝업 & 59 & 0.741 & \textbf{0.551} \\
펀딩 & 80 & 0.426 & 0.370 \\
도서 & 163 & 0.612 & 0.271 \\
웹툰 & 711 & 0.485 & 0.139 \\
\hline
\end{tabular}
\end{center}

아이돌의 ``천장 0.797''은 거의 전부 외운 것이다. 판 수준에서 섞음이
\textbf{0.2080} 이다.

\section{정직한 눈금은 $+$0.013 이다}

\begin{figure}[t]
\centering
\includegraphics[width=\columnwidth]{figs/ceiling.pdf}
\caption{도메인마다 섞음(과적합 바닥) $\to$ 속천장 구간, 그 안에 지금(파랑)과
유보 안 CV(초록). 초록이 파랑 왼쪽이면 여지가 음수다.}
\end{figure}

과적합을 안 타는 눈금은 유보 \emph{안}에서의 5겹 CV 다 --- 그 도메인의
유보 라벨로 배우되 자기가 안 본 행에서 잰다. 배포 모형에게는 불공평하게
유리한 조건이다(유보 라벨을 본다).

\begin{center}\small
\begin{tabular}{lrrrr}
\hline
도메인 & $n$ & 지금 & CV & 여지 \\
\hline
아이돌 & 25 & 0.154 & 0.406 & $+$0.252 \\
\textbf{도서} & 163 & 0.321 & 0.498 & $\mathbf{+0.177}$ \\
웹툰 & 711 & 0.381 & 0.431 & $+$0.051 \\
팝업 & 59 & 0.391 & 0.423 & $+$0.033 \\
애니 & 606 & 0.487 & 0.473 & $-$0.014 \\
모바일 & 441 & 0.537 & 0.512 & $-$0.026 \\
세계애니 & 300 & 0.520 & 0.487 & $-$0.033 \\
게임 & 180 & 0.640 & 0.601 & $-$0.039 \\
펀딩 & 80 & 0.236 & 0.190 & $-$0.046 \\
\hline
\textbf{판} & & \textbf{0.4569} & \textbf{0.4697} & $\mathbf{+0.0128}$ \\
\hline
\end{tabular}
\end{center}

\begin{figure}[t]
\centering
\includegraphics[width=\columnwidth]{figs/board.pdf}
\caption{판 수준. 섞음과 속천장 사이 어디에 있느냐가 아니라, 지금과 CV 의
간격이 남은 여지다.}
\end{figure}

\textbf{남은 여지가 $+$0.0128 이다.} 노트 245가 그은 선(0.010)을 겨우 넘는다.
그리고 \textbf{아홉 중 다섯이 음수}다 --- 풀링 모형이 그 도메인의 유보 라벨을
직접 보고 배운 모형보다 낫다. 노트 126이 ``도메인 고유 유연성은 비싸다''고
쓴 것의 가장 강한 형태다. \textbf{풀링은 타협이 아니라 이득이다.}

\section{도서 하나가 남았다}

여지가 크고 표본도 믿을 만한 곳은 도서뿐이다($+$0.177, $n{=}163$, 몫 6.2\%
--- 다 메우면 판 $+$0.011). 아이돌 $+$0.252 는 $n{=}25$ 에 섞음이 0.68 이라
숫자로 쓸 수 없다(노트 248).

왜인지는 \textbf{라벨을 안 보고} 보인다.

\begin{center}\small
\begin{tabular}{lrrr}
\hline
도메인 & 학습 & 유보 & 비 \\
\hline
\textbf{도서} & \textbf{80} & \textbf{163} & \textbf{0.49} \\
게임 & 259 & 180 & 1.44 \\
아이돌 & 54 & 25 & 2.16 \\
애니 & 1,467 & 606 & 2.42 \\
웹툰 & 2,106 & 711 & 2.96 \\
모바일 & 1,559 & 441 & 3.54 \\
펀딩 & 320 & 80 & 4.00 \\
세계애니 & 2,648 & 300 & 8.83 \\
\hline
\end{tabular}
\end{center}

\textbf{도서는 학습이 유보보다 적은 유일한 도메인이다}(팝업은 16 대 59 로
더 심하지만 몫이 2.3\%다). 게다가 그 80건이 \textbf{1999--2025} 에 흩어져
있고, 축 넷(\texttt{media\_push} · 위키 셋)이 덮음 0\% 다. 도서의 여지는
모형이 아니라 \textbf{수집}이다.

비로 여지를 예측할 수 있으면 규약이 된다. 쟀더니 spearman $-$0.52(n$=$9,
믿을 만한 일곱만 보면 $-$0.54)다. 방향은 맞지만 \textbf{n 이 아홉이라 법칙이라
부를 수 없다} --- 도서가 이상치인 것이지 비가 여지를 설명하는 것이 아니다.
적어 두되 규약으로 올리지 않는다.

\section{프로그램 수준의 결론}

이 노트가 말하는 것은 한 실험의 결과가 아니라 \textbf{어디까지 왔는가}다.

\begin{itemize}
\item 축 열일곱으로 남은 여지는 $+$0.013 이다. 노트 240--251이 판을 못
움직인 것은 방법이 나빠서가 아니라 \textbf{움직일 것이 거의 없어서}다.
\item 다섯 도메인은 이미 천장 위에 있다. 그곳에서 모형을 손보는 것은
\emph{원리적으로} 손해다.
\item 남은 것은 도서 $+$0.011 과 웹툰 $+$0.014(몫 27.3\% $\times$ 0.051)뿐이고,
도서는 수집 문제다.
\item \textbf{다음 이득은 더 나은 모형이 아니라 새 정보에서 온다.} 노트 222의
도메인 천장, 노트 220의 무리 천장이 같은 말을 각각의 층에서 했고, 이제
전체 층에서 확인됐다.
\end{itemize}

\paragraph{남은 것}
\begin{center}\small
\begin{tabular}{ll}
\hline
갈래 & 왜 \\
\hline
도서 수집 & 학습 80건. 늘리면 $+$0.011 까지 \\
웹툰 $+$0.051 & 몫이 27\%라 판 $+$0.014. 무엇이 모자라나 \\
새 축 & 여지가 $+$0.013 이면 새 정보 말고 길이 없다 \\
팝업 16건 & 학습이 유보의 4분의 1. 몫이 작아 판엔 안 보인다 \\
\hline
\end{tabular}
\end{center}

\paragraph{규약} 모형을 손보기 전에 \textbf{유보 안 CV 로 여지를 먼저 잰다.}
음수면 그 도메인에서는 무엇을 해도 손해다. 노트 220 · 222의 천장 계산에
한 줄 더 붙는 셈이고, 셋 다 ``재기 전에 알 수 있는 것''이다.

\end{document}
